% !TeX encoding = UTF-8
\documentclass[variant=docente,cover=institutional,mode=screen]{ua-engineering-v6-article}
% !TeX encoding = UTF-8
\UASetUniversity{Universidad Austral}
\UASetFaculty{Facultad de Ingeniería}
\UASetDepartment{Departamento de Computación}
\UASetCampus{Pilar}
\UASetCareer{Ingeniería Informática}
\UASetCommission{Comisión A}
\UASetProfessor{Equipo docente}
\UASetSemester{Segundo cuatrimestre 2026}
\UASetCourse{Sistemas Distribuidos}
\UASetDate{2026}


\UASetClassNumber{12}
\UASetClassTitle{Chaos Engineering y validación bajo falla}
\UASetDocKind{Guía docente}

\title{Clase 12 --- Guía docente}
\author{\UAProfessor}
\date{\UADate}

\begin{document}
\maketitle

\section{Objetivo pedagógico profundo}

La meta de esta clase es transformar la idea de resiliencia desde una propiedad “supuestamente presente” hacia una propiedad experimentalmente validada.

La idea que debe quedar instalada es:

\begin{center}
\textbf{un sistema resiliente no se presume: se verifica bajo perturbaciones controladas}
\end{center}

\section{Resultados de aprendizaje esperados}

Al finalizar, el estudiante debería poder:

\begin{itemize}
\item explicar qué diferencia a Chaos Engineering de testing clásico;
\item formular una hipótesis experimental razonable;
\item identificar steady state y métricas observables;
\item diseñar una perturbación controlada con blast radius acotado;
\item interpretar resultados sin caer en lecturas superficiales.
\end{itemize}

\section{Estructura sugerida para 90 minutos}

\subsection*{Bloque 1 (15 min): por qué testing clásico no alcanza}
Arrancar con un contraejemplo:
\begin{itemize}
\item sistema pasa todos los tests;
\item en producción, una latencia alta en una dependencia dispara retries y colapsa el sistema.
\end{itemize}

\paragraph{Objetivo didáctico}
Mostrar que la validación distribuida requiere modelos de falla más realistas.

\subsection*{Bloque 2 (20 min): Chaos Engineering}
Definir la disciplina, insistiendo en:
\begin{itemize}
\item hipótesis;
\item steady state;
\item perturbación;
\item medición;
\item aprendizaje.
\end{itemize}

\paragraph{Objetivo didáctico}
Evitar que el estudiante lo entienda como simple vandalismo controlado.

\subsection*{Bloque 3 (20 min): fault injection}
Presentar tipos de fallas:
\begin{itemize}
\item crash;
\item latencia;
\item pérdida;
\item saturación;
\item partición.
\end{itemize}

Discutir qué propiedad diferente valida cada una.

\paragraph{Objetivo didáctico}
Que el alumno vea que no todas las fallas testean lo mismo.

\subsection*{Bloque 4 (15 min): blast radius y seguridad}
Explicar:
\begin{itemize}
\item por qué los experimentos deben ser graduales;
\item cómo se limita el radio de impacto;
\item cuándo abortar.
\end{itemize}

\paragraph{Objetivo didáctico}
Introducir responsabilidad operacional como parte del diseño experimental.

\subsection*{Bloque 5 (20 min): ejemplo integrado}
Trabajar un caso concreto:
\begin{itemize}
\item Payments lento;
\item hipótesis sobre breaker y retries;
\item métricas de steady state;
\item interpretación de resultados.
\end{itemize}

\paragraph{Objetivo didáctico}
Conectar resiliencia, observabilidad y validación empírica en un solo caso.

\section{Qué explicar y por qué}

\subsection*{1. Testing clásico}
Porque da el contraste necesario para entender qué agrega Chaos Engineering.

\subsection*{2. Steady state}
Porque sin steady state el experimento no tiene criterio de éxito o fracaso.

\subsection*{3. Fault injection}
Porque es el mecanismo que hace visible el comportamiento bajo perturbación.

\subsection*{4. Blast radius}
Porque el curso debe formar criterio técnico y también responsabilidad de operación.

\subsection*{5. Integración con observabilidad}
Porque un experimento sin observabilidad produce ruido, no conocimiento.

\section{Preguntas disparadoras recomendadas}

\begin{itemize}
\item ¿Qué hipótesis concreta querés validar?
\item ¿Qué variable observás para afirmar que el sistema “sigue bien”?
\item ¿Qué parte del sistema vas a perturbar y por qué?
\item ¿Qué aprendés si el sistema falla? ¿Y si no falla?
\item ¿Cómo sabés cuándo abortar el experimento?
\end{itemize}

\section{Errores frecuentes del alumnado}

\begin{itemize}
\item creer que Chaos Engineering es “tirar abajo cosas”;
\item proponer experimentos sin steady state;
\item diseñar perturbaciones sin límites de impacto;
\item olvidar rollback o abort;
\item no distinguir entre síntoma medido y propiedad validada.
\end{itemize}

\section{Ejercicio en vivo sugerido}

Proponer un sistema de checkout con:
\begin{itemize}
\item API;
\item Auth;
\item Payments;
\item Orders.
\end{itemize}

Pedir al curso que diseñe:
\begin{itemize}
\item una hipótesis sobre latencia alta en Payments;
\item el steady state relevante;
\item la perturbación;
\item el blast radius;
\item los criterios para abortar.
\end{itemize}

\paragraph{Objetivo}
Que el estudiante vea el experimento como una práctica disciplinada y no como una ocurrencia.

\section{Momento crítico de la clase}

El punto más importante aparece cuando el estudiante entiende que:

\begin{center}
\textbf{la pregunta útil no es “si el sistema falla”, sino “cómo falla y qué evidencia tengo sobre ese comportamiento”}
\end{center}

Ese punto debe remarcarse explícitamente.

\section{Criterio de éxito docente}

La clase salió bien si el estudiante puede:

\begin{itemize}
\item formular una hipótesis experimental concreta;
\item definir steady state medible;
\item proponer una perturbación controlada;
\item limitar el blast radius;
\item interpretar el resultado como conocimiento operacional.
\end{itemize}

\section{Conexión con la clase siguiente}

Esta clase deja abierta una transición natural hacia:
\begin{itemize}
\item confiabilidad;
\item SLO/SLA/SLI;
\item error budgets;
\item gestión de operación basada en objetivos.
\end{itemize}

\begin{uakeybox}
Pregunta puente: una vez que podemos observar y experimentar bajo falla, ¿cómo definimos formalmente qué significa que el sistema sea confiable para sus usuarios?
\end{uakeybox}


\section{Plan de conducción ampliado}
La clase debe alternar explicación, predicción y contraste. Antes de revelar cada mecanismo, pedir al grupo que anticipe qué puede salir mal; después, volver al mismo escenario y preguntar qué propiedad mejora y qué costo aparece. Cada bloque conceptual debe cerrar con una decisión de diseño observable.
\subsection*{Secuencia recomendada}
\begin{enumerate}
\item Recuperar el prerequisito de la clase anterior.
\item Presentar un caso mínimo que falle y solicitar hipótesis.
\item Formalizar el concepto central de Chaos Engineering y validación bajo falla, distinguiendo seguridad, progreso y supuestos.
\item Resolver un ejemplo completo en pizarra, incluyendo una falla intermedia.
\item Comparar una alternativa de diseño y explicitar trade-offs.
\item Cerrar con una pregunta del Trabajo Final y una pregunta de defensa oral.
\end{enumerate}
\section{Diagnóstico de comprensión durante la clase}
No preguntar solamente ``¿se entendió?''. Usar preguntas discriminantes: ``¿qué cambia si el mensaje llega dos veces?'', ``¿qué propiedad sigue siendo cierta si cae un nodo?'', ``¿esto demuestra seguridad o progreso?'', ``¿qué observa un cliente externo?''. Si el estudiante responde solo con nombres de patrones, pedir una ejecución concreta.
\section{Criterios para corregir ejercicios}
Una respuesta excelente declara supuestos, construye una secuencia verificable, nombra la garantía con precisión, reconoce un costo y conecta la decisión con el dominio. Penalizar afirmaciones absolutas sin condiciones.
\section{Vinculación con el Trabajo Final Integrador}
Reservar entre 5 y 10 minutos para que cada grupo registre una decisión con: problema, alternativa elegida, alternativa descartada y razón. Esto crea trazabilidad entre las clases y las entregas acumulativas.
\section{Lectura docente previa}
Revisar la bibliografía indicada en los apuntes y preparar al menos un contraejemplo que no aparezca literalmente en las diapositivas. El objetivo es poder responder preguntas de frontera sin convertir la clase en una lectura del deck.

\section{Arquitectura didáctica v9 para 180 minutos}
\subsection*{Bloque 1 - desafío inicial (20 min)}
Presentar una situación donde aparezca perturbación sin hipótesis, blast radius excesivo, señal insuficiente y falso positivo. Pedir predicciones individuales antes de introducir vocabulario. El objetivo es hacer visibles los supuestos intuitivos y recuperar el prerequisito de la clase anterior.

\subsection*{Bloque 2 - construcción conceptual (35 min)}
Formalizar testing distribuido, fault injection, steady state, hipótesis, blast radius y diseño experimental. Separar seguridad, progreso y experiencia del usuario. Explicar no solo \emph{qué} hace cada mecanismo, sino \emph{por qué} existe y qué supuesto permite que funcione.

\subsection*{Bloque 3 - modelo y figura (35 min)}
Construir en pizarra una línea temporal, grafo, log, máquina de estados o tabla de decisiones. Recorrer una ejecución nominal y una adversarial. La representación debe quedar como referencia para los apuntes y el ejercicio principal.

\subsection*{Pausa (10 min)}
Recoger una pregunta anónima o un punto de confusión. Elegir una pregunta que permita distinguir síntoma, causa y propiedad.

\subsection*{Bloque 4 - caso trabajado con falla (40 min)}
Aplicar el mecanismo al caso conductor. Introducir una falla a mitad de ejecución y detenerse en cada transición: quién sabe qué, qué se persiste, qué garantía sigue vigente y qué observa el cliente.

\subsection*{Bloque 5 - taller de decisión (25 min)}
Los grupos aplican P-F-G-S-T a su Trabajo Final. Deben producir un ADR breve con alternativa elegida, alternativa descartada y evidencia pendiente. La evidencia de esta clase es baseline, perturbación, métricas, criterio de aborto y conclusión reproducible.

\subsection*{Bloque 6 - cierre y exit ticket (15 min)}
Cada estudiante responde: propiedad principal, supuesto decisivo, costo aceptado y condición bajo la cual cambiaría de diseño. Usar las respuestas para iniciar la clase siguiente.

\section{Qué explicar, por qué y cómo verificarlo}
\begin{itemize}
\item Explicar el problema antes del producto, porque reconocer la familia de problema es el resultado cognitivo central.
\item Explicar el invariante, porque permite evaluar configuraciones y no solo repetir un flujo nominal.
\item Explicar el límite, porque perturbación sin hipótesis, blast radius excesivo, señal insuficiente y falso positivo puede invalidar supuestos o reducir progreso.
\item Verificar comprensión mediante baseline, perturbación, métricas, criterio de aborto y conclusión reproducible, no mediante una definición aislada.
\end{itemize}

\section{Preguntas socráticas}
\begin{itemize}
\item ¿Qué sabe cada actor en este punto de la ejecución?
\item ¿Qué propiedad sigue siendo cierta si la respuesta nunca llega?
\item ¿Qué parte es safety y cuál es liveness?
\item ¿Qué alternativa más simple resolvería el problema si relajamos una garantía?
\item ¿Qué observación refutaría nuestra hipótesis de diseño?
\end{itemize}

\section{Errores previsibles y estrategia de corrección}
\begin{itemize}
\item Si el alumno responde con un nombre de tecnología, pedir actores, mensajes y estados.
\item Si confunde timeout con falla, construir dos ejecuciones indistinguibles.
\item Si promete todas las garantías, pedir comportamiento bajo partición o saturación.
\item Si mide solo rendimiento, preguntar cómo evidencia la propiedad semántica.
\item Si la solución es excesiva, pedir la alternativa mínima y el costo marginal de la garantía fuerte.
\end{itemize}

\section{Criterios de evaluación formativa}
Una respuesta excelente declara supuestos, recorre una ejecución, nombra con precisión la garantía, reconoce el costo y propone evidencia reproducible. Una respuesta suficiente identifica el mecanismo y el trade-off principal. Una respuesta insuficiente usa slogans, omite fallas o confunde producto con propiedad.

\section{Preparación docente y bibliografía}
Lectura previa recomendada: Principles of Chaos Engineering; Jepsen; literatura de fault injection y testing distribuido. Preparar un contraejemplo adicional y una variante del caso en la que cambie el costo del error; esto evita convertir la clase en una lectura de diapositivas.

\end{document}
